% Exercise 1.1: paper reviews.
% Add the group members' names to \author before submission.
\documentclass[11pt,a4paper]{article}
\usepackage[T1]{fontenc}
\usepackage[utf8]{inputenc}
\usepackage[english]{babel}
\usepackage{lmodern}
\usepackage[margin=2.2cm]{geometry}
\usepackage{microtype}
\usepackage[hidelinks]{hyperref}
\setlength{\parindent}{0pt}
\setlength{\parskip}{0.6em}
\setlength{\emergencystretch}{2em}
\clubpenalty=10000
\widowpenalty=10000
\title{Exercise 1.1: Paper Reviews}
\author{}
\date{}

\begin{document}
\maketitle
\vspace{-2em}

\section*{Evolution of the Graphics Processing Unit (GPU)}
{\small W. J. Dally, S. W. Keckler, and D. B. Kirk.
\emph{IEEE Micro}, vol. 41, no. 6, pp. 42--51, 2021.}

\subsection*{1. Primary contribution}
This article provides a historical synthesis of how GPUs evolved from
fixed-function graphics hardware into programmable platforms for graphics,
high-performance computing, and deep learning.
The authors connect major developments, including programmable shaders,
CUDA-enabled computation, and Tensor Cores, to the requirements of these
applications.
By examining hardware alongside programming tools and libraries, they
explain how architectural capabilities and a supporting software ecosystem
jointly enabled the adoption of GPUs beyond graphics.

\subsection*{2. Key insight of the contribution}
The key insight is that GPU progress is driven by feedback between applications
and architecture: programmability attracts new workloads, whose requirements
then motivate specialized hardware.
Tensor Cores illustrate this process by accelerating matrix
multiply-and-accumulate operations and amortizing instruction overhead
across many arithmetic operations.
Together with suitable reduced-precision formats and software support,
such specialization explains why the inference-performance trend called
Huang's law exceeds the gains attributable to transistor scaling alone.

\subsection*{3. Opinion and reaction}
I found it interesting that the drive for better video games helped create
the computing platform behind the modern AI boom.
The article presents this progression logically and makes the industry's
direction understandable without relying on excessively detailed technical
explanations.
I also appreciated its emphasis on software, which makes clear why
computational power alone is not enough to make a platform useful.

\newpage
\section*{A Bridging Model for Parallel Computation}
{\small Leslie G. Valiant.
\emph{Communications of the ACM}, vol. 33, no. 8, pp. 103--111, 1990.}

\subsection*{1. Primary contribution}
Valiant introduces the bulk-synchronous parallel (BSP) model as a common
framework between parallel software and hardware, with a role similar to
that of the von Neumann model in sequential computing.
The model combines processing and memory components, point-to-point
communication, and controlled synchronization in units called supersteps.
The paper supports this proposal with results on implementing shared-memory
features and parallel algorithms on BSP, and on implementing BSP efficiently
using different communication networks.

\subsection*{2. Key insight of the contribution}
The key insight is that a program should expose more parallel work than
the physical processors can execute at once.
This parallel slackness gives the compiler room to distribute work,
pipeline communication, and absorb uneven memory-access loads.
Together with synchronization at a suitable granularity, it allows many
low-level details to be managed automatically with only a constant-factor
efficiency loss, provided enough parallelism and communication capacity
are available.

\subsection*{3. Opinion and reaction}
I found it interesting that the paper treats a shared model for software
and hardware as an essential part of making parallel computing practical.
The idea of exposing more parallelism than the machine can execute at once
seems counterintuitive at first, but its role in handling communication
and uneven workloads makes the argument convincing.
I would have liked a more detailed programming example, since the efficiency
arguments are harder to connect to an implementation than the basic model
itself.

\end{document}
